\documentclass[../../main.tex]{subfiles}

\graphicspath{{../../image/}{../../figure/}}

\begin{document}

\section{阻尼与松弛的作用}\label{section:ch05-阻尼与松弛的作用}

前几节在 Cauchy 问题的经典解框架下展示了凸熵的稳定性作用。本节将表明，即使没有凸熵，稳定性也可由阻尼或松弛产生，而当这些机制与双曲性相平衡时，它们可产生整体光滑解，即使初始振幅并不很小。

我们从研究一类双曲平衡律方程组开始，它可写成\eqref{align:ch04-黏性组}的形式，
即状态向量分裂为两部分 $U=(V,W)$，通量 $G=(F,H)$，源
\begin{align}
Q(U)=\binom{0}{R(V,W)},\label{align:ch05-Q-split}
\end{align}
使得
\begin{align}
R(V,0)=0,\qquad V\in\mathscr{B}_{\rho_{0}}.\label{align:ch05-R-equilibrium}
\end{align}
于是 $W=0$ 是平衡流形。守恒律方程组为
\begin{align}
\partial_{t}V+\sum_{\alpha=1}^{m}\partial_{\alpha}F_{\alpha}(V,W)&=0,\label{align:ch05-split1}\\
\partial_{t}W+\sum_{\alpha=1}^{m}\partial_{\alpha}H_{\alpha}(V,W)=R(V,W).\label{align:ch05-split2}
\end{align}
我们假设 $R(V,0)=0$ 是 $R(V,W)=0$ 在 $\overline{\mathscr{B}_{\rho}}$ 上的唯一解，特别地
\begin{align}
\mathrm{D}_{W}R(V,0)\ \text{在}\ \overline{\mathscr{B}_{\rho_{0}}}\ \text{上正定},\label{align:ch05-DWR-positive}
\end{align}
从而平衡流形是渐近稳定的。

装备有熵 $\eta(U)$ 的方程组\eqref{align:ch05-split1}、\eqref{align:ch05-split2}在\cref{section:ch04-消失黏性方法}中已被考虑过。
那里证明了能量耗散
\begin{align}
\int_{\mathbb{R}^{m}}\mathrm{D}_{W}\eta(V,0)R(V,W)\,\mathrm{d}x\leqslant 0,\label{align:ch05-dissipation}
\end{align}
对 $W$ 的高阶取负号，正是子特征条件\eqref{align:ch04-耗散条件}。
在连续介质物理中出现的许多方程组里，$\mathrm{D}_{W}\eta(V,0)R(V,W)$ 实际上正比于 $|W|^{2}$。典型情形是 $W$ 为速度偏差而 $R$ 给出摩擦力。

我们在此研究\eqref{align:ch05-split1}、\eqref{align:ch05-split2} Cauchy 问题经典解的存在性。本节的第一个结果是\cref{theorem:ch05-凸熵局部存在}的推广：

\begin{theorem}\label{theorem:ch05-damping-existence}
假设方程组\eqref{align:ch05-split1}、\eqref{align:ch05-split2}在平衡流形 $W=0$ 附近双曲，且对任意固定 $V\in\overline{\mathscr{B}_{\rho_{0}}}$ 关于 $W$ 严格双曲。设初值 $(V_{0},W_{0})\in H_{\ell}$，某 $\ell>\frac{m}{2}+1$，取值于 $\overline{\mathscr{B}_{\rho}}$。则存在 $T_{\infty}\leqslant\infty$ 及 $\mathbb{R}^{m}\times[0,T_{\infty})$ 上唯一连续可微、取值于 $\overline{\mathscr{B}_{\rho}}$ 的函数 $(V,W)$，它是 Cauchy 问题\eqref{align:ch05-split1}、\eqref{align:ch05-split2}在时间区间 $[0,T_{\infty})$ 上的经典解。此外，$(V,W)(\cdot,t)\in\bigcap_{k=0}^{\ell}C^{k}([0,T_{\infty});H_{\ell-k})$。区间 $[0,T_{\infty})$ 是极大的，即若 $T_{\infty}<\infty$，则
\begin{align}
\int_{0}^{T_{\infty}}\|\|\nabla(V,W)(\cdot,t)\|\|_{L^{\infty}}\,\mathrm{d}t=\infty,\label{align:ch05-damping-blowup1}
\end{align}
且/或
\begin{align}
\limsup_{t\to T_{\infty}}\|(V,W)(\cdot,t)\|_{L^{\infty}}=\rho.\label{align:ch05-damping-blowup2}
\end{align}
\end{theorem}

证明沿用\cref{theorem:ch05-凸熵局部存在}的路线，借助于粘性消去法。
然而，由于此处熵不必凸，对称化子须在 $W$ 上局部构造，由 $\mathrm{D}_{W}^{2}\eta(V,0)$ 诱导，并借\eqref{align:ch05-DWR-positive}在 $W=0$ 附近正定。

下一个命题处理小振幅初值的整体存在性。

\begin{theorem}\label{theorem:ch05-damping-global}
在\cref{theorem:ch05-damping-existence}的假设下，假设此外方程组\eqref{align:ch05-split1}、\eqref{align:ch05-split2}装备有熵 $\eta(V,W)$，它在平衡流形 $W=0$ 附近正定，且满足耗散条件\eqref{align:ch05-dissipation}。则存在 $\delta>0$，使得任意初值 $(V_{0},W_{0})\in H_{\ell}$、$\ell>\frac{m}{2}+1$、取值于 $\overline{\mathscr{B}_{\rho}}$ 且满足 $\|(V_{0},W_{0})\|_{\ell}<\delta$，生成定义在整个 $\mathbb{R}^{m}\times[0,\infty)$ 上的唯一经典解 $(V,W)$。此外，
\begin{align}
\|W(\cdot,t)\|_{0}\leqslant c\|(V_{0},W_{0})\|_{0}e^{-\kappa t},\qquad t\geqslant0,\label{align:ch05-W-decay}
\end{align}
对某 $\kappa>0$，而
\begin{align}
\|V(\cdot,t)\|_{0}\leqslant c\|(V_{0},W_{0})\|_{0},\qquad t\geqslant0.\label{align:ch05-V-bounded}
\end{align}
\end{theorem}

证明通过能量估计进行：以对称矩阵 $A(U)=\mathrm{D}^{2}\eta(U)$ 乘方程组，分部积分并利用 $\mathrm{D}_{W}\eta R$ 在 $W=0$ 附近的负定性。

我们接下来研究松弛诱导阻尼的典型情形。考虑形式
\begin{align}
\partial_{t}U(x,t)+\sum_{\alpha=1}^{m}\partial_{\alpha}G_{\alpha}(U(x,t))=\frac{1}{\mu}Q(U(x,t)),\label{align:ch05-relaxed-system}
\end{align}
其中 $\mu>0$ 是小的松弛参数。我们假设 $Q(0)=0$，$\mathrm{D}Q(0)$ 的核张成由状态向量分裂 $U=(V,W)$ 所对应的子空间，且 $Q$ 在 $W$ 方向耗散。形式展开渐近导出\emph{松弛系统}
\begin{align}
\partial_{t}V(x,t)+\sum_{\alpha=1}^{m}\partial_{\alpha}\hat{G}_{\alpha}(V(x,t))=0,\label{align:ch05-relaxed-hyperbolic}
\end{align}
其中 $\hat{G}_{\alpha}(V)=G_{\alpha}(V,\hat{W}(V))$，$\hat{W}(V)$ 由 $Q(V,\hat{W})=0$ 隐式确定。松弛系统是双曲的，且原系统\eqref{align:ch05-relaxed-system}的解在 $\mu\to0$ 时收敛到\eqref{align:ch05-relaxed-hyperbolic}的解。

为使小振幅解整体存在，原系统须满足\emph{子特征条件}
\begin{align}
\mathrm{D}_{W}H(V,\hat{W}(V))[\mathrm{D}_{W}R(V,\hat{W}(V))]^{-1}\mathrm{D}_{W}H(V,\hat{W}(V))^{\top}\geqslant0,\label{align:ch05-subcharacteristic}
\end{align}
它刻画了松弛波速不超过松弛系统特征波速。

作为一个简单的例子，考虑松弛方程组
\begin{align}
\partial_{t}u+\partial_{x}f(u,w)&=0,\label{align:ch05-example1}\\
\partial_{t}w+\partial_{x}h(u,w)=\frac{1}{\mu}[a(u)w-f^{\prime}(u)w],\label{align:ch05-example2}
\end{align}
其松弛系统为标量守恒律
\begin{align}
\partial_{t}u(x,t)+\partial_{x}f(u(x,t))=0,\label{align:ch05-relaxed-scalar}
\end{align}
特征速度 $\hat{\lambda}=f^{\prime}(u)$。如前所述，子特征条件
\begin{align}
-a(u)\leqslant f^{\prime}(u)\leqslant a(u)\label{align:ch05-subchar-scalar}
\end{align}
是使系统\eqref{align:ch05-example1}、\eqref{align:ch05-example2}存在使源耗散的凸熵 $\eta(u,v)$ 的必要前提。我们将在后续章节回到系统\eqref{align:ch05-example1}、\eqref{align:ch05-example2}，特别地，将看到\eqref{align:ch05-subchar-scalar}以严格不等式形式成立也是凸熵存在的充分条件。

在\cref{section:ch04-消失黏性方法}中，我们遇到了诱导熵产生的另一机制，即粘性。相关系统呈\eqref{align:ch04-黏性组}形式，

熵产生编码在\eqref{align:ch04-黏性熵恒等式}右端第二项中。
类比上述关于带源源系统引入的术语，当与 $\mathrm{D}^{2}\eta B_{\alpha\beta}$ 相联系的二次型半正定时，称粘性项是\emph{耗散}的；特别地，当该二次型正定或耗散半正定且\eqref{align:ch04-耗散条件}以 $a>0$ 成立时，称为\emph{耗散定}的。
耗散定粘性使系统\eqref{align:ch04-黏性组}抛物化，
此时 Cauchy 问题适定并容许很光滑的解。然而，连续介质物理中出现的系统里常见的粘性项仅为耗散半正定。事实上，此类系统常呈形式
\begin{align}
\begin{cases}
\partial_{t}V+\sum_{\alpha=1}^{m}\partial_{\alpha}F_{\alpha}(V,W)=0,\\
\partial_{t}W+\sum_{\alpha=1}^{m}\partial_{\alpha}H_{\alpha}(V,W)=\mu\sum_{\alpha,\beta=1}^{m}\partial_{\alpha}[B_{\alpha\beta}(V,W)\partial_{\beta}W],
\end{cases}\label{align:ch05-viscous-system}
\end{align}
与\eqref{align:ch04-黏性组}对照。例子见\eqref{align:ch04-黏性组}。


在各种设定下，带半正定耗散粘性的系统 Cauchy 问题解的存在性与长时间行为已被广泛研究。新兴理论在范围、方法论与结论上与上述带耗散半正定源的平衡律方程组 Cauchy 问题理论紧密平行。特别地，对适定性，当粘性仅为耗散半正定时，需要补充条件：由\eqref{align:ch04-黏性组}关于平衡态 $U=0$ 线性化所得的系统
\begin{align}
\partial_{t}V+\sum_{\alpha=1}^{m}\mathrm{D}G_{\alpha}(0)\partial_{\alpha}V=\mu\sum_{\alpha,\beta=1}^{m}B_{\alpha\beta}(0)\partial_{\alpha}\partial_{\beta}V\label{align:ch05-linearized-viscous}
\end{align}
不容许无阻尼行波解。当\emph{Kawashima 条件}
\begin{align}
\sum_{\alpha,\beta=1}^{m}v_{\alpha}v_{\beta}B_{\alpha\beta}(0)R_{i}(v;0)\neq0,\qquad v\in\mathbb{S}^{m-1},\ i=1,\ldots,n\label{align:ch05-kawashima}
\end{align}
成立时，此要求被满足。这是带耗散半正定源的平衡律方程组的 Kawashima 条件的对应。

若将\eqref{align:ch05-split1}、\eqref{align:ch05-split2}的解按小松弛参数 $\mu$ 的幂作形式展开，即可揭示粘性松弛与由刚性源诱导的松弛之间的紧密联系。在零阶 $\mu^{0}$，得到松弛双曲系统\eqref{align:ch05-relaxed-hyperbolic}。在下一阶 $\mu^{1}$，结果是带粘性的系统：
\begin{align}
\partial_{t}V+\sum_{\alpha=1}^{m}\partial_{\alpha}\hat{G}(V)=\mu\sum_{\alpha,\beta=1}^{m}\partial_{\alpha}[\hat{B}_{\alpha\beta}(V)\partial_{\beta}V].\label{align:ch05-viscous-relaxed}
\end{align}
此外，粘性关于\eqref{align:ch05-relaxed-hyperbolic}的熵 $\hat{\eta}(V)$ 是耗散的。子特征条件\eqref{align:ch05-subcharacteristic}是稳定性的另一相关表现。

为避免对一般系统确定\eqref{align:ch05-viscous-relaxed}中系数 $\hat{B}_{\alpha\beta}$ 的繁琐计算，我们在简单系统\eqref{align:ch05-example1}、\eqref{align:ch05-example2}的背景下阐明上述内容。我们已看到标量守恒律\eqref{align:ch05-relaxed-scalar}是松弛系统。为达到下一个 $\mu$ 阶，将 $v=f(u)+\mu w$ 代入\eqref{align:ch05-example2}，借助\eqref{align:ch05-example1}$_1$ 消去 $\partial_{t}u$，再舍去所有 $\mu$ 阶项，得
\begin{align}
w=[f^{\prime}(u)^{2}-a(u)^{2}]\partial_{x}u.\label{align:ch05-w-substitute}
\end{align}
最后将 $v=f(u)+\mu w$ 代入\eqref{align:ch05-example1}$_1$，得
\begin{align}
\partial_{t}u+\partial_{x}f(u)=\mu\partial_{x}\{[a(u)^{2}-f^{\prime}(u)^{2}]\partial_{x}u\},\label{align:ch05-final-viscous}
\end{align}
即\eqref{align:ch05-linearized-viscous}的形式。子特征条件\eqref{align:ch05-subchar-scalar}与\eqref{align:ch05-final-viscous}中粘性耗散性之间的关系现已清楚。

在连续介质物理中，人们遇到大量演化系统，其特征为非线性对流诱导的波放大与某种耗散共存并竞争：前者在远离平衡时占优，后者在平衡附近占优，从而保证大范围光滑解的存在。此类系统一般用本节采用的类似方法处理，即"能量"型估计，揭示放大与阻尼之间的平衡。这一主题已拥有大量文献，超出本书范围。尽管如此，为让读者对具此类特征的系统的广泛多样性有一印象，下面记录几个代表性例子，少量相关文献列于\cref{section:ch05-注释}。

我们从所谓 Euler--Poisson 系统开始：
\begin{align}
\begin{cases}
\partial_{t}\rho+\operatorname{div}(\rho v^{\top})=0,\\
\partial_{t}(\rho v)+\operatorname{div}(\rho vv^{\top})+\operatorname{grad}p(\rho)=a\rho\,\operatorname{grad}\psi,\\
\Delta\psi=b(\rho-\bar{\rho}).
\end{cases}\label{align:ch05-euler-poisson}
\end{align}
当 $a=-1$、$b=4\pi G$、$\bar{\rho}=0$ 时，\eqref{align:ch05-euler-poisson}描述气体在自身质量产生的引力场中的流动，即"吸引"情形。在相反的"排斥"情形，即 $a,b$ 均为正常数时，该系统模拟等离子体中电子的运动。就此而言，电子集合被视为密度 $\rho$、压力 $p(\rho)$ 的弹性流体，以速度 $v$ 流动；而重得多的离子假设静止，仅提供正比于 $\bar{\rho}$ 的均匀正电荷背景。电子与离子的组合电荷正比于 $\rho-\bar{\rho}$，产生静电势 $\psi$，从而产生使电子运动的电场 $\operatorname{grad}\psi$。如\eqref{align:ch05-split1}、\eqref{align:ch05-split2}，此处处理的是 Euler 方程双曲系统，带有由\eqref{align:ch05-euler-poisson}$_3$ Poisson 方程导出的反馈机制诱导的源，且至少在电子流动无旋（$\operatorname{curl}v=0$）时是耗散的。由\cref{section:ch03-例子}回忆，无旋流动只要光滑就保持无旋。
已证明：充分光滑、无旋、接近平衡态 $\rho=\bar{\rho}$、$v=0$、$\psi=0$ 的 Cauchy 初值生成整体定义的光滑解。另一方面，从远离平衡态出发的解一般在有限时间内发展出奇性。

下列半导体流体力学模型系统的情形类似：
\begin{align}
\begin{cases}
\partial_{t}\rho+\operatorname{div}(\rho v^{\top})=0,\\
\partial_{t}(\rho v)+\operatorname{div}(\rho vv^{\top})+\operatorname{grad}p(\rho)+\mu^{-1}\rho v=a\rho\,\operatorname{grad}\psi,\\
\Delta\psi=b(\rho-\bar{\rho}),
\end{cases}\label{align:ch05-semiconductor}
\end{align}
其中 $\mu>0$ 是松弛参数。注意\eqref{align:ch05-semiconductor}结合了\eqref{align:ch05-split1}、\eqref{align:ch05-split2}与\eqref{align:ch05-euler-poisson}中遇到的耗散机制。

带内摩擦（如粘性与热传导）连续介质的平衡律产生类似行为的系统。原因在于可将这些介质的谱系追溯到弹性，从而尽管所得系统不双曲，却继承双曲性的特征，产生与内摩擦诱导的阻尼相竞争的失稳波放大机制。

第一个例子是\eqref{align:ch04-NS}，
它描述带 Newton 粘性的热导热粘弹性流体的流动。内摩擦显现在第二个与第三个方程的右端，而第一个方程保持其双曲特征。

另一具类似特征的例子是系统
\begin{align}
\begin{cases}
\partial_{t}u-\partial_{x}v=0,\\
\partial_{t}v-\partial_{x}\sigma(u,\theta)=0,\\
\partial_{t}[\varepsilon(u,\theta)+\frac{1}{2}v^{2}]-\partial_{x}[\sigma(u,\theta)v]=\partial_{x}q(u,\theta,\partial_{x}\theta),
\end{cases}\label{align:ch05-thermoelastic}
\end{align}
它描述 Lagrangian 坐标下热导热弹性介质的直线运动。其中 $u$ 是应变（变形梯度），$v$ 是速度，$\theta$ 是（绝对）温度，$\sigma$ 是应力，$\varepsilon$ 是内能，$q$ 是热流，参考密度取为 1。为与\eqref{align:ch02-2-5-28}与\eqref{align:ch02-2-5-29}一致，

材料响应函数 $\varepsilon,\sigma,q$ 须满足 $\varepsilon_{u}=\sigma-\theta\sigma_{\theta}$ 与 $q(u,\theta,g)g\geqslant0$。这些应补充以自然假设 $\sigma_{u}>0$、$\varepsilon_{\theta}>0$ 与 $q_{g}>0$。此处内摩擦由热扩散提供。

另一性质相似的内摩擦由粘弹性连续介质中的\emph{fading memory}（衰减记忆）诱导，其中质点 $x$ 在时刻 $t$ 的应力 $\sigma$ 不再像弹性材料那样仅由 $(x,t)$ 处的变形梯度决定，而还依赖于 $x$ 处变形梯度的过去历史。平衡律于是以泛函-偏微分方程表达。捕捉记忆阻尼效应的一个简单一维模型系统为
\begin{align}
\begin{cases}
\partial_{t}u(x,t)+\partial_{x}v(x,t)=0,\\
\partial_{t}v(x,t)+\partial_{x}p(u(x,t))+\int_{-\infty}^{t}k^{\prime}(t-\tau)\partial_{x}q(u(x,\tau))\,\mathrm{d}\tau=0,
\end{cases}\label{align:ch05-fading-memory}
\end{align}
其中 $k$ 是 $[0,\infty)$ 上光滑可积的松弛核，$k(\tau)>0$、$k^{\prime}(\tau)<0$、$k^{\prime\prime}(\tau)\geqslant0$（$\tau\in[0,\infty)$），且 $p^{\prime}(u)>k(0)q^{\prime}(u)>0$。注意\eqref{align:ch05-fading-memory}与\eqref{align:ch05-example1}、\eqref{align:ch05-example2}紧密相关，后者当 $f\equiv0$ 时可改写为
\begin{align}
\begin{cases}
\partial_{t}u(x,t)+\partial_{x}v(x,t)=0,\\
\partial_{t}v(x,t)+\partial_{x}p(u(x,t))+\int_{-\infty}^{t}\left[\exp\left(-\frac{t-\tau}{\mu}\right)\right]^{\prime}\partial_{x}p(u(x,\tau))\,\mathrm{d}\tau=0.
\end{cases}\label{align:ch05-fading-memory2}
\end{align}

上述系统\eqref{align:ch05-thermoelastic}与\eqref{align:ch05-fading-memory}的共同点是：接近平衡态的光滑初值生成整体定义的光滑解，而从"大"初值出发的光滑解一般在有限时间内爆破。见\cref{section:ch05-注释}的相关文献。

作用于连续介质物理平衡律方程组的另一衰减机制是色散。当空间维数大且解保持接近平衡时，色散特别有效。如\cref{section:ch04-Euler方程}所见，
在完全非线性系统（如 Euler 方程）中，色散可延迟但不能阻止波破裂。然而，在满足所谓\emph{null condition}（零条件）的较温和非线性系统中，色散使接近平衡态初值的 Cauchy 问题整体光滑解的存在性成为可能。作为典型例子，考虑等熵弹性动力学方程组\eqref{align:ch03-等熵热弹性系统}。
为方便起见，假设参考空间与物理空间重合，参考位形 $F\equiv I$ 是各向同性平衡态，从而内能 $\varepsilon(F)$ 是右伸长张量主不变量 $(J_{1},J_{2},J_{3})$ 的函数\eqref{align:ch02-2-5-21}。
进一步假设 $\varepsilon(F)$ 秩一凸并满足零条件
\begin{align}
\sum_{i,j,k=1}^{3}\sum_{\alpha,\beta,\gamma=1}^{3}\frac{\partial^{3}\varepsilon(F)}{\partial F_{i\alpha}\partial F_{j\beta}\partial F_{k\gamma}}v_{\alpha}v_{\beta}v_{\gamma}v_{i}v_{j}v_{k}=0,\label{align:ch05-null-condition}
\end{align}
在 $F=I$ 处对任意 $v\in\mathbb{R}^{3}$ 成立。则接近 $(I,0)$ 的初值 $(F_{0},v_{0})$ 在适当 Sobolev 空间中，Cauchy 问题容许唯一整体定义的经典解。对各向同性不可压缩弹性介质，相关零条件自动满足。关于这些问题的文献汗牛充栋，少量样本引于\cref{section:ch05-注释}。

\end{document}
